\documentclass[../../main.tex]{subfiles}
\begin{document}
% SOURCE p=251,252; printed=236-237; object=ch08-schwartz-definition; region=bottom-upper
\begin{definition}[Schwartz 空间]\label{definition:ch08-schwartz-space}
另一个需要的 $C^\infty$ 函数空间是 \textbf{Schwartz 空间} $\mathcal S$：其中的函数及其所有阶导数在无穷远处均比 $|x|$ 的任意负幂衰减得更快。准确地说，对非负整数 $N$ 及多重指标 $\alpha$，定义
\begin{align*}
\|f\|_{(N,\alpha)}&=\sup_{x\in\mathbb R^n}(1+|x|)^N|\partial^\alpha f(x)|,\\
\mathcal S&=\{f\in C^\infty:\|f\|_{(N,\alpha)}<\infty\text{ 对所有 }N,\alpha\}.
\end{align*}
\end{definition}
% SUPPLEMENT: id=v14-ch08-schwartz-definition-details; parent=ch08-schwartz-definition; kind=omitted-proof+motivation+hypotheses; origin=independent-derivation
\begin{remark}
\textbf{动机与逐条件说明（重构新增）。}这个空间同时控制任意阶求导和任意多项式增长，为 Fourier 分析中的反复分部积分提供自然的函数类。要求\emph{所有} $N$，而非某一个固定衰减阶；要求\emph{所有} $\alpha$，而非只控制函数本身，因为求导后的函数也要保留快速衰减。权 $(1+|x|)^N$ 在原点不为零，既控制局部大小，也控制无穷远。定义确实给出“快于每个幂”的极限：用 $N+1$ 阶权的有界性可得 $(1+|x|)^N|\partial^\alpha f(x)|\leqslant C/(1+|x|)\to0$。

每个 $\|\cdot\|_{(N,\alpha)}$ 在 $\mathcal S$ 上实际上是范数。齐次性和三角不等式直接成立。若 $\partial^\alpha f=0$，当 $\alpha_j>0$ 时，令 $g=\partial^{\alpha-e_j}f\in\mathcal S$；$\partial_jg=0$ 使 $g$ 在每条平行第 $j$ 坐标轴的直线上为常数，而 $g$ 沿该直线趋于无穷远时趋零，故 $g=0$。有限次降低阶数最终得到 $f=0$；$\alpha=0$ 时立即成立。这里导数仍在 $\mathcal S$，是因为定义中的全部权和全部导数估计对它同样适用。
\end{remark}
% END SUPPLEMENT: id=v14-ch08-schwartz-definition-details
% SOURCE p=252; printed=237; object=ch08-schwartz-examples; region=upper
\begin{example}[Schwartz 函数的基本例子]\label{example:ch08-schwartz-examples}
容易找到 $\mathcal S$ 中的函数，例如任意多重指标 $\alpha$ 给出的 $f_\alpha(x)=x^\alpha e^{-|x|^2}$。此外，$C_c^\infty\subseteq\mathcal S$。
\end{example}
% SUPPLEMENT: id=v14-ch08-schwartz-examples-proof; parent=ch08-schwartz-examples; kind=omitted-proof+motivation+hypotheses; origin=independent-derivation
\begin{remark}
\textbf{完整补证（重构新增）。}由乘积法则，每次求导都使“多项式乘 $e^{-|x|^2}$”仍有这种形式，故 $\partial^\beta f_\alpha=P_{\alpha,\beta}(x)e^{-|x|^2}$。若多项式次数为 $d$，则其绝对值不超过 $C(1+|x|)^d$；指数衰减使 $(1+r)^{N+d}e^{-r^2}$ 在 $r\geqslant0$ 上有界。例如由 $e^s\geqslant s^k/k!$，选 $2k>N+d$ 即得无穷远处趋零，紧区间上连续有界。因此全部 Schwartz 范数有限。对紧支撑光滑函数，每阶导数的支撑都包含于原支撑，且在这个紧集上连续有界；权也有界，故同样属于 $\mathcal S$。这两类例子分别展示无紧支撑的快速衰减与紧支撑两种来源。
\end{remark}
% END SUPPLEMENT: id=v14-ch08-schwartz-examples-proof
% SOURCE p=252; printed=237; object=ch08-schwartz-lp; region=upper-middle
\begin{proposition}[Schwartz 函数各阶导数的可积性]\label{proposition:ch08-schwartz-lp}
若 $f\in\mathcal S$，则每个多重指标 $\alpha$ 和每个 $p\in[1,\infty]$ 都满足 $\partial^\alpha f\in L^p$。
\end{proposition}
\begin{proof}
对所有 $N$ 均有 $|\partial^\alpha f(x)|\leqslant C_N(1+|x|)^{-N}$。当 $p<\infty$ 时，取 $N>n/p$，由\cref{corollary:ch02-power-integrability}，$(1+|x|)^{-N}\in L^p$，因而结论成立。$p=\infty$ 则直接使用 $N=0$ 的估计。
\end{proof}
% SUPPLEMENT: id=v14-ch08-schwartz-lp-details; parent=ch08-schwartz-lp; kind=motivation+hypotheses; origin=independent-derivation
\begin{remark}
\textbf{动机与条件（重构新增）。}在单位球上，该权函数有界；在球外其 $p$ 次方至多 $|x|^{-Np}$，径向积分为常数倍的 $\int_1^\infty r^{n-1-Np}\,\mathrm dr$，恰在 $Np>n$ 时有限。严格不等号保证幂指数小于 $-1$；边界值对应的比较积分发散。Schwartz 条件允许自由选取足够大的 $N$，所以每个有限 $p$ 都可处理；$p=\infty$ 只需有界性。这里每一阶导数都享有同样衰减，因此结论覆盖全部 $\alpha$，为后续积分及求导交换准备了可积控制。
\end{remark}
% END SUPPLEMENT: id=v14-ch08-schwartz-lp-details
% SOURCE p=252; printed=237; object=ch08-schwartz-frechet; number=8.2; region=middle
\begin{proposition}\label{proposition:ch08-schwartz-frechet}
$\mathcal S$ 配备由范数 $\|\cdot\|_{(N,\alpha)}$ 生成的拓扑，是 Fréchet 空间。
\end{proposition}
\begin{proof}
主要需要证明完备性。若 $\{f_k\}$ 是 $\mathcal S$ 中的 Cauchy 序列，则对全部 $N,\alpha$ 都有 $\|f_j-f_k\|_{(N,\alpha)}\to0$。特别地，对每个 $\alpha$，$\partial^\alpha f_k$ 一致收敛到某个连续函数 $g_\alpha$。记 $e_j$ 为第 $j$ 个标准单位向量。由微积分基本定理，
\begin{align*}
f_k(x+te_j)-f_k(x)=\int_0^t\partial_jf_k(x+se_j)\,\mathrm ds.
\end{align*}
令 $k\to\infty$，得到
\begin{align*}
g_0(x+te_j)-g_0(x)=\int_0^t g_{e_j}(x+se_j)\,\mathrm ds.
\end{align*}
因此 $g_{e_j}=\partial_jg_0$。将同样论证用于每个 $\partial^\alpha f_k$，再按 $|\alpha|$ 归纳，得到 $g_\alpha=\partial^\alpha g_0$，因而 $g_0\in C^\infty$。最后，由加权范数下的 Cauchy 条件可得 $\|f_k-g_0\|_{(N,\alpha)}\to0$，故 $g_0\in\mathcal S$，序列在 $\mathcal S$ 中收敛。
\end{proof}
% SUPPLEMENT: id=v14-ch08-schwartz-completeness-details; parent=ch08-schwartz-frechet; kind=omitted-proof+motivation+hypotheses; origin=independent-derivation
\begin{remark}
\textbf{拓扑及极限步骤的完整补证（重构新增）。}向量空间性由各加权估计的三角不等式给出。这是可数个范数组成的族，把它们枚举为 $q_1,q_2,\ldots$，则
\begin{align*}
d(f,g)=\sum_{j=1}^\infty2^{-j}\min\{1,q_j(f-g)\}
\end{align*}
是平移不变度量：三角不等式由 $\min(1,a+b)\leqslant\min(1,a)+\min(1,b)$ 得到，正定性由范数正定性得到。有限头部逐项趋零、尾部几何级数任意小，说明所有 $q_j$ 收敛蕴含度量收敛；反向由任一固定项 $2^{-j}\min(1,q_j)$ 不超过 $d$ 得到。同样论证适用于 Cauchy 条件。有限个范数开球的交是凸邻域基，故该拓扑局部凸且 Hausdorff；这也与\cref{theorem:ch05-seminorm-topology,proposition:ch05-seminorm-metrization}一致。

原证中每个 $\partial^\alpha f_k$ 在一致范数下 Cauchy，逐点取复数的完备极限，再由一致 Cauchy 条件得到一致收敛，故 $g_\alpha$ 连续且有界。对任意固定 $x,t,j,\alpha$，积分差的绝对值至多为
$|t|\|\partial^{\alpha+e_j}f_k-g_{\alpha+e_j}\|_u\to0$，所以
\begin{align*}
g_\alpha(x+te_j)-g_\alpha(x)=\int_0^t g_{\alpha+e_j}(x+se_j)\,\mathrm ds.
\end{align*}
连续性及基本定理给 $\partial_jg_\alpha=g_{\alpha+e_j}$，这逐阶证明了所有极限确实是同一个函数的导数。

再固定 $N,\alpha,\varepsilon>0$，取 $k_0$ 使 $j,k\geqslant k_0$ 时 $\|f_j-f_k\|_{(N,\alpha)}<\varepsilon$。固定 $k\geqslant k_0$ 和 $x$，令 $j\to\infty$，得到
$(1+|x|)^N|g_\alpha(x)-\partial^\alpha f_k(x)|\leqslant\varepsilon$。对 $x$ 取上确界后，$\|g_0-f_k\|_{(N,\alpha)}\leqslant\varepsilon$。取一个这样的 $k$ 还给 $\|g_0\|_{(N,\alpha)}\leqslant\|f_k\|_{(N,\alpha)}+\varepsilon<\infty$。这同时证明极限的快速衰减和全部加权范数收敛；仅有无权的一致收敛不足以替代这一步。由于 $\varepsilon$ 任意，度量完备，故为 Fréchet 空间。可数性用于可度量性，全部阶数与全部权的 Cauchy 条件用于使极限留在原空间；这就是该拓扑的动机。
\end{remark}
% END SUPPLEMENT: id=v14-ch08-schwartz-completeness-details
% SOURCE p=252; printed=237; object=ch08-schwartz-equivalences; number=8.3; region=lower
\begin{proposition}\label{proposition:ch08-schwartz-equivalences}
若 $f\in C^\infty$，则以下条件等价：$f\in\mathcal S$；对全部多重指标 $\alpha,\beta$，$x^\beta\partial^\alpha f$ 有界；对全部多重指标 $\alpha,\beta$，$\partial^\alpha(x^\beta f)$ 有界。
\end{proposition}
\begin{proof}
当 $|\beta|\leqslant N$ 时，$|x^\beta|\leqslant(1+|x|)^N$。反向估计中先设 $N\geqslant1$。单位球面上连续函数 $\sum_{j=1}^n|x_j|^N$ 严格为正，故有正最小值 $\delta$；两边的 $N$ 次齐次性给出 $\sum_j|x_j|^N\geqslant\delta|x|^N$。又 $\delta\leqslant1$，因为可代入 $e_1$，所以
\begin{align*}
(1+|x|)^N
&\leqslant2^N(1+|x|^N)
\leqslant2^N\left[1+\delta^{-1}\sum_{j=1}^n|x_j|^N\right]\\
&\leqslant2^N\delta^{-1}\sum_{|\beta|\leqslant N}|x^\beta|.
\end{align*}
这证明前两个条件等价；$N=0$ 时直接用 $\beta=0$ 即可。最后，乘积法则使每个 $\partial^\alpha(x^\beta f)$ 为有限多个 $x^\delta\partial^\gamma f$ 的线性组合，反向也成立，见\cref{exercise:ch08-leibniz}。因此后两个条件等价。
\end{proof}
% SUPPLEMENT: id=v14-ch08-schwartz-equivalence-details; parent=ch08-schwartz-equivalences; kind=omitted-proof+motivation+hypotheses; origin=independent-derivation
\begin{remark}
\textbf{有限估计、反向展开及零阶边界（重构新增）。}对于固定的 $N,\alpha$，上式只涉及有限多个 $\beta$，因此各 $x^\beta\partial^\alpha f$ 的有界性确实给出一个有限的统一上界。另一方面，\cref{exercise:ch08-leibniz}逐阶解出 $x^\beta\partial^\alpha f$：把 $\partial^\alpha(x^\beta f)$ 展开后，除主项外各项的导数阶数都更低；按 $|\alpha|$ 归纳，就可把这些低阶项改写成有限多个 $\partial^\gamma(x^\delta f)$。因而反向并不是未经证明的交换求导与乘法。

原书在球面估计中未单列 $N=0$。此时 $\delta=n$，若 $n>1$，显示的不等式链最后一步不成立；这里将那条链明确用于 $N\geqslant1$，零阶直接处理，结论不变。动机是用单项式权取代径向权，便于与微分交换。无限光滑保证两种排列的所有阶导数均存在；全部 $\alpha,\beta$ 是整体 Schwartz 性的要求，一个有限子族只能给有限阶控制。
\end{remark}
% END SUPPLEMENT: id=v14-ch08-schwartz-equivalence-details
\subfile{section-01-translations.tex}
\end{document}
